%% notes-tex/figure-3-gate/sections/5-status.tex
%% SOURCE: hand-authored status, evening of 2026-10-05. Round numbers are those of
%% experiments/019-simb-multimodal/results/joint_checkpoint_readout.json (v19 tests, the
%% v20_partial block); benchmark numbers are from the W&B projects torchcell_019_small_bench
%% and torchcell_019_grid_bench and the task logs of GilaHyper jobs 3267, 3271, 3273 and 3274;
%% queue states from squeue on IGB and Delta at the times given (CDT). Each is a reading of a
%% log or a results file, not a result, unless the paragraph says otherwise.

\clearpage
\section{Status on the evening of 2026-10-05\secstatus{todo}}
\label{sec:status}

What has finished, what is running, and what each framing rests on tonight. The mockups of
\cref{sec:mockups} were redrawn from the result files as they stand at 22:30; the only panel
whose data changed since the morning is \cref{fig:mockup-a}d.

\pfstep{Round v19 is complete and null on both registered tests} All twelve split seeds
finished training on IGB; the readout on file reads eleven, split seed 11 having been
unfinished when it was taken. Expression at the registered window (epochs 1{,}000 to
1{,}200), joint minus single head: mean $-0.001$, standard deviation across partitions
$0.030$, 5 of 11 above zero, one-sided $t$ test $p = 0.55$. Proteome at epochs 200 to 400:
mean $-0.017$, 1 of 11 above zero, $p = 0.64$ against the non-inferiority margin of
$-0.015$. Joint training on identical both-label rows buys nothing on either head at this
budget. \cref{fig:mockup-b}c and \cref{fig:mockup-si-1}e draw these eleven; the twelfth is
added when the readout is rerun.
\src{joint_checkpoint_readout.json, v19.tests}

\pfstep{Round v20 is complete on two of twelve partitions} Tasks 0 and 1 of IGB job 2423190
reached epoch 1{,}199 after 31 hours each on \file{cabbi}; tasks 2 and 3 were at epochs 36
to 88 at 22:25 and tasks 4 to 11 wait under a two-at-a-time throttle. On the two complete
partitions, each statistic the mean validation Pearson per feature over the latest 20
epochs of the conditioned run against the same epochs of the unconditioned v19 arm:

\begin{center}\footnotesize
\begin{tabular}{@{}llrrrr@{}}
\toprule
head & given & split seed & conditioned & unconditioned & permuted control \\
\midrule
proteome   & expression & 0 & 0.168 & 0.112 & 0.126 \\
proteome   & expression & 1 & 0.095 & 0.016 & 0.029 \\
expression & proteome   & 0 & 0.125 & 0.117 & 0.122 \\
expression & proteome   & 1 & 0.126 & 0.083 & 0.003 \\
\bottomrule
\end{tabular}
\end{center}

The proteome head conditioned on the strain's own expression sits $+0.056$ and $+0.079$
above the unconditioned arm, and the permuted-expression control gains only $+0.014$ and
$+0.013$, so the gain uses the strain's own measurement and not the presence of a second
input. The expression head conditioned on the proteome gains $+0.008$ and $+0.043$, and on
split seed 0 its permuted control gains as much ($+0.005$), so there is no separation there
yet. The bar the plan set is the per-partition ridge from genotype plus the measured
modality, $0.236$ on expression and $0.214$ on the proteome on one split; both conditioned
arms are below it on both partitions. Two partitions are two draws from a distribution whose
between-partition spread is $0.02$ to $0.03$ (\cref{sec:review}); this is a reading, not a
result. W\&B runs \texttt{7tiaq7pe}, \texttt{fg63qx22} (conditioned proteome),
\texttt{3o61ceay}, \texttt{izzszopm} (its permuted control), \texttt{7a7tc6xb},
\texttt{qvmr4wyl} (conditioned expression), \texttt{zxlow2i4}, \texttt{u6p781n5}, project
\texttt{torchcell\_019\_prot\_v20}.
\src{joint_checkpoint_readout.json, v20_partial.runs}

\pfstep{Track A: the matched baselines have run and are unread} IGB job 2423302, the linear
and nearest-neighbor baselines on exactly the strains of the both-label store at twelve
split seeds, finished on 2026-10-04 and has not been read. Until it is, \cref{fig:mockup-a}f
and \cref{fig:mockup-b}d keep the caveat that baselines and model share a seed number and
not a store.

\pfstep{Depth and width do not set the epoch time} A three-hour window on the four GilaHyper
cards (RTX 6000 Ada, 44.4 GiB) was spent timing the model at reduced size rather than
training it, since convergence cannot be read in three hours. Four cells at zero
data-loader workers, four to eight runs sharing a card: the v19 model on the 3{,}328-d
stack (6.68 million parameters), ProtT5 alone at the same trunk (1.45 million), ProtT5 at
three layers (1.15 million), and ProtT5 at width 54 with two layers (0.74 million). Card
memory was 8 to 10 GB per run at every shape, so eight runs on a card ran out of memory
(three of eight), six sat at the limit, and four is the ceiling; the memory is the
activations over 6{,}600 gene tokens at batch 32, not the parameters. Steady-state epoch
time at four per card was 78 s for the full model and 80 to 81 s for every ProtT5 shape. The
same held with the cards fed: with three persistent workers per run the epoch fell to 35 to
40 s at four per card and did not move with depth. Delta's own v11 round logged 81 to 134 s
per epoch at one run per A40 with zero workers. So at both worker settings the trunk is not
where the time goes. The runs are in W\&B project \texttt{torchcell\_019\_small\_bench}
(GilaHyper jobs 3267, 3271 and 3273) and \texttt{torchcell\_019\_grid\_bench} (job 3274).
\src{notes/experiments.019-simb-multimodal.scripts.gh_small_model_bench.md}

\begin{center}\footnotesize
\begin{tabular}{@{}lrrr@{}}
\toprule
cell (four runs per card) & workers per run & seconds per epoch per run & card memory \\
\midrule
3{,}328-d stack, width 90, six layers & 0 & 78 & 36 GB \\
ProtT5, width 90, six layers & 0 & 80 to 81 & 36 GB \\
ProtT5, width 90, three layers (six per card) & 0 & 80 & at the limit \\
ProtT5, width 54, two layers (eight per card) & 0 & 81 & out of memory \\
ProtT5, width 90, six layers & 3 & 35 to 40 & 38 GB \\
ProtT5, width 90, six layers & 1 (eight CPUs) & 58 to 64 & 38 GB \\
ProtT5, width 90, two layers & 3 & 37 to 40 & 33 GB \\
ProtT5, width 90, three layers & 3 & 38 to 41 & 37 GB \\
ProtT5, width 90, four layers & 3 & 39 to 43 & 36 GB \\
ProtT5, width 90, six layers & 3 & 39 to 42 & 37 GB \\
\bottomrule
\end{tabular}
\end{center}

The last four rows are the depth-by-operator grid (job 3274) at epoch 9, each row the range
over its four operators; the operators sit within 3 s of each other at every depth. The
persistent-worker setting is the one lever found: the Delta rule of zero workers dates from
July, when workers were re-spawned every epoch and re-imported torch from Lustre; persistent
workers pay that import once. Hypothesis (untested): the time scales with batch size times
gene count, through the node tensors, the per-gene readout and loss, the metric passes and
the eval-mode passes, and one place it can go is the perturbation operator, which the code
runs as a Python loop over the strains in a batch
(\file{torchcell/models/equivariant_cell_graph_transformer.py}), one cross-attention
over all genes per strain; a one-step profile would settle it. The cached-encoder speedup of
\cref{sec:review} was measured on a CPU and does not transfer to a GPU on this evidence.

\pfstep{Decision taken: the small trunk is ProtT5 alone at width 90 and six layers} Since
depth and width buy no time and no packing, the v21 configuration
(\file{conf/cgt_expr_v21_small.yaml}) changes one thing against v19, the input: ProtT5
alone, 1.45 million parameters against 6.68, with every other setting inherited. Hypothesis
(untested until v21 reads out): the smaller trunk reaches the v19 single-head level on these
rows.

\pfstep{Round v21 is on Delta} Track B, wave 1, submitted at 21:58 on
\file{bfjt-delta-gpu} at source 9ca207582: nine arms on twelve split seeds, 108 runs, as 27
independent single-A40 jobs of four runs each (jobs 22690625 to 22690651), one persistent
worker per run on eight CPUs, 1{,}200 epochs, a two-day limit, no dependency between jobs.
A canary of four reference copies at three epochs with a test pass (job 22690624) runs
beside them; if it fails, the round is canceled. The queue was faster than the scheduler's
estimate: \texttt{sbatch --test-only} placed the first pack at 16:08 the next day, and
fifteen packs were running by 22:50, the canary 54 minutes in and still importing the store
from Lustre. The arms, each one change against the reference S\_ref (the v21 configuration as
is): S\_mask, the masked objective restored; S\_sink, the sigmoid-gated operator with the
gate open; S\_prop2, the deletion propagated two hops along the graphs; S\_nodrop, operator
attention dropout off; S\_stack, the four-part embedding stack in place of ProtT5 (the size
control); S\_prot, the single-head proteome model on this trunk; S\_basis64, the rank-64
response basis; S\_hadam, the elementwise product at the operator. Scoring is
pre-registered in the configuration header: the mean of the validation Pearson per feature
over epochs 1{,}000 to 1{,}200 (the proteome arm over 200 to 400); an arm is carried to wave
2 if it beats S\_ref on at least nine of twelve split seeds; a run whose prediction spread
ratio is below $0.05$ at epoch 200 is relaunched with a new initialization seed. The store
is v19's: fig3\_proteome restricted to the 1{,}349 genotypes carrying both labels, synced to
Delta tonight.
\src{experiments/019-simb-multimodal/scripts/delta_expr_v21_small.slurm,
conf/cgt_expr_v21_small.yaml, scripts/gh_expr_008_arm.sh}

\pfstep{The operator question has three candidates under test} \cref{fig:mockup-si-2}e
still waits on the author's choice between Fig.~1f and the Methods. Three operators that
give each gene its own response at one deletion are in v21 at twelve paired seeds each
(S\_sink, S\_hadam, S\_basis64) and in the GilaHyper grid tonight on split seed 0 at four
depths. The grid's operator read, validation Pearson and its spread at about epoch 200 by
operator within each depth, lands when job 3274 reaches its wall at 00:00 (its two-hour
limit was cut to 96 minutes to fit the window, so the runs end near epoch 160 rather than
200) and is appended to this document in its next version. Sixteen runs on one split seed cannot decide
the operator; they size what v21 decides.

\pfstep{What this leaves for the two framings} Framing A rests on two complete partitions of
v20 and the ridge triangle; the conditioned proteome arm clears its permuted control by
$0.04$ to $0.07$ on both and is below the ridge bar on both. Framing B is unchanged: v19 is
null, the matched baselines are on disk unread, and no trained model on file beats the
nearest-neighbor key on matched strains. The first number that can move either framing is
the v21 readout on 2026-10-07 and 2026-10-08, and the next partitions of v20 as \file{cabbi}
releases them.
